\scanpage{75}
% lines 1–2, heading 4 and upper-left note
\textcircled{4} Relation de $\pi_i,T_i$
\[
\boxed{\pi_i(x)=\unsure{\chi|\gamma}x-T_i}
\]
% slanted
\margin{\unsure{remonte|remarque} aux généralités sur les $\chi$-trivialisations, n'a rien à voir avec l'op. $*$}
% lines 3–9, heading 5 and check on a submodule
\textcircled{5} Description de la $\chi$-trivialisation de $E$
\[
\boxed{\pi_E(x_1*x_2)=\pi_1(x_1)\otimes\pi_2(x_2)}
\]

Si $x_2=u_2\in\Gamma L_2$, on a $\inserted{x\overset{\mathrm{def}}{=}x_1*x_2={}}x_1*u_2=\pi_1(x_1)\otimes u_2\deleted{\unsure{\in E}}\inserted{\in\Gamma L_1\otimes L_2}$ et \deleted{on aura} le premier mb doit donc être égal à $\chi x=\chi\,\pi_1(x_1)\otimes u_2$, le deuxième est $\pi_1(x_1)\otimes\chi u_2$ puisque $\pi_2(u_2)=\chi u_2$, et c'est OK. \unsure{Item|Idem} pour $x_1=u_1$, $x_2$ quelcon.
% lines 10–15, corollary
Cor. On aura, pour $x_1\in\Gamma P_1$, $x_2\in\Gamma P_2$
\[
\underset{\substack{\shortparallel\\\chi x_1\otimes x_2-T}}{\pi_E(x_1*x_2)}
=\underbrace{(\chi x_1-T_1)}_{b_1}\otimes\underbrace{(\chi x_2-T_2)}_{b_2}
\]
donc
\[
\boxed{T=\chi(\,x_1\mathbin{\unsure{*|\otimes}}x_2)-b_1\otimes b_2}
\]
où
\[
b_i=\pi_{E^i}(x_i)=\chi x_i-T_i
\]
% lines 16–21, long slanted note at right
% slanted
\margin{NB La donnée de $x_1,x_2$ d'où $x_1*x_2$ définit \unsure{des} scindages\\
$E_1\simeq L_1\oplus\underline O\quad(u_1,\lambda_1)\mapsto\lambda_1 x_1+u_1$\\
$E_2\simeq L_2\oplus\underline O\quad(u_2,\lambda_2)\to\lambda_2 x_2+u_2$\\
$E\simeq\underbrace{L_1\otimes L_2}_{L}\oplus\unsure{U|O},\quad(u,\lambda)\mathbin{-}\lambda x_1*x_2+\deleted{\unsure{x}}u$\\
On aura alors\\
$(u_1,\lambda_1)\deleted{\unsure{\otimes}}*(u_2,\lambda_2)=$\\
$(\lambda_1\deleted{\unsure{b_2}}b_1\otimes u_2+u_1\otimes\lambda_2 b_2+\chi u_1\otimes u_2,\lambda_1\lambda_2)$}
% lines 22–26, verification at foot
Vérifions
\[
\boxed{\chi T=\underset{\substack{\shortparallel\\\chi x_1-b_1}}{T_1}*
\underset{\substack{\shortparallel\\\chi x_2-b_2}}{T_2}}
\]
le $2^{\mathrm{\unsure{d|o|e}}}$ \unsure{mb} est
\[
\begin{aligned}
&\chi^2 x_1*x_2-\chi\,\underbrace{\pi_1(x_1)}_{(\chi x_1-T_1)=b_1}\otimes b_2
-\chi\,b_1\otimes\underbrace{\pi_2(x_2)}_{(\chi x_2-T_2)=b_2}
+\underbrace{b_1*b_2}_{\chi b_1\otimes b_2}\\
&=\chi^2 x_1*x_2-\chi\,b_1\otimes b_2=\chi(\overbrace{\chi x_1*x_2-b_1\otimes b_2}^{T})
\end{aligned}
\]
c'est donc bien \deleted{ce qui avait \unsure{précédé|par $\chi$ de}} $\chi T$, OK
% archive sheet number: 74
